\section{Eisenstein Series and Quantum Computation}


\subsection{Overview of Eisenstein Series}

Eisenstein series are fundamental objects in the theory of modular
forms. For a complex variable $z$ in the upper half-plane
$\mathcal{H} = \{z \in \mathbb{C} : \operatorname{Im}(z) > 0\}$
and an even integer weight $k \geq 4$, the Eisenstein series is:

\begin{equation}\label{eq:eisenstein}
G_k(z) = \sum_{\substack{(c,d) \in \mathbb{Z}^2 \\ (c,d) \neq (0,0)}}
\frac{1}{(cz + d)^k}
\end{equation}

The sum ranges over the integer lattice $\mathbb{Z}^2$ excluding the
origin. The series converges absolutely for $\operatorname{Re}(k) > 2$
and transforms as a modular form of weight $k$ for
$\mathrm{SL}_2(\mathbb{Z})$:

\[
G_k\!\left(\frac{az+b}{cz+d}\right) = (cz+d)^k \, G_k(z),
\qquad
\begin{pmatrix} a & b \\ c & d \end{pmatrix} \in \mathrm{SL}_2(\mathbb{Z})
\]

In the context of this work, we consider Eisenstein series associated
to cusps of the Hilbert--Blumenthal modular group, which provide
additional structure for multi-dimensional lattice summations.

\subsection{Quantum Representation of Complex Numbers}\label{sec:qrep}

To evaluate Eisenstein series on a quantum computer, we must encode
complex arithmetic in the qubit model. Let $z = a + bi$ where
$a, b \in \mathbb{R}$. We use a fixed-point representation with
$p$ bits of precision for each component:

\begin{definition}[Quantum Complex Register]
A quantum complex number on $2p$ qubits is a state
$\ket{z} = \ket{a_1 a_2 \ldots a_p} \otimes \ket{b_1 b_2 \ldots b_p}$
where $a_i, b_i \in \{0,1\}$ encode the fixed-point binary
representations of $\operatorname{Re}(z)$ and $\operatorname{Im}(z)$.
\end{definition}

Arithmetic operations (addition, multiplication, reciprocal) on
quantum complex registers require $O(p \log p)$ gates via quantum
Fourier transform techniques~\cite{nielsen2010}. The key operation
for Eisenstein evaluation is the computation of $(cz+d)^{-k}$, which
decomposes into:

\begin{enumerate}[nosep]
  \item Multiply: $w \leftarrow cz + d$ \quad ($O(p \log p)$ gates)
  \item Power: $w^k$ via repeated squaring \quad ($O(k \cdot p \log p)$ gates)
  \item Invert: $w^{-k}$ via Newton iteration \quad ($O(p^2)$ gates)
\end{enumerate}

\subsection{Quantum State Construction}

Each term $(c,d)$ in the Eisenstein series is represented by a
computational basis state $\ket{c,d}$ with $c,d$ encoded as signed
integers on $\ell$ qubits each. We prepare the uniform superposition
over a truncated lattice $\Lambda_N = \{(c,d) \in \mathbb{Z}^2 :
|c|, |d| \leq N, (c,d) \neq (0,0)\}$:

\begin{equation}\label{eq:superposition}
\ket{\Psi_0} = \frac{1}{\sqrt{|\Lambda_N|}}
\sum_{(c,d) \in \Lambda_N} \ket{c,d} \otimes \ket{0}
\end{equation}

where $|\Lambda_N| = (2N+1)^2 - 1$. This superposition is prepared
in $O(\ell)$ Hadamard gates followed by a single ancilla-controlled
rejection of the $(0,0)$ state.

\subsection{Quantum Summation via State Interference}

The core of the algorithm applies a unitary $U_z$ that computes the
term value in an ancilla register:

\[
U_z \ket{c,d}\ket{0} = \ket{c,d}\ket{(cz+d)^{-k}}
\]

The Eisenstein sum is then encoded in the amplitude of a target state.
Define the projection operator $P_{\mathrm{sum}}$ that extracts the
sum of ancilla values. By linearity of quantum mechanics:

\begin{equation}\label{eq:interference}
U_z \ket{\Psi_0} = \frac{1}{\sqrt{|\Lambda_N|}}
\sum_{(c,d) \in \Lambda_N} \ket{c,d}\ket{(cz+d)^{-k}}
\end{equation}

The sum $\tilde{G}_k^{(N)}(z) = \sum_{(c,d) \in \Lambda_N} (cz+d)^{-k}$
is not directly observable as an amplitude (amplitudes are normalized),
but it can be estimated via quantum phase estimation on a
controlled-$U_z$ circuit. The estimation requires
$O(\sqrt{|\Lambda_N|} / \epsilon)$ queries for additive error
$\epsilon$, a quadratic improvement over classical summation.

%% ============================================================
